% ===========================================================================
\appendix

\section{Example LLM Adjudication}
\label{sec:appendix-example}

This appendix presents a complete example of LLM adjudication from Tournament 9, Round 10.
The dictator is the \textbf{Mirror} strategy and the recipient is the \textbf{Universalizability} strategy.
This example illustrates the recursive reasoning required when two sophisticated conditional strategies interact.

\subsection{Prompt Structure}

The LLM receives a structured prompt containing:
\begin{enumerate}
  \item Game rules and instructions
  \item The dictator's strategy (in full)
  \item The recipient's strategy (in full)
  \item Instructions to apply the dictator's strategy and output a JSON decision
\end{enumerate}

\subsection{Complete Prompt}

\begin{small}
\begin{verbatim}
You are an expert at determining how a decision procedure should be applied.

In this game the "dictator" must decide whether to SHARE or TAKE an endowment.
- **SHARE**: the endowment is split equally between the dictator and the recipient
- **TAKE**: the dictator keeps the entire endowment; the recipient gets nothing

You must follow only the dictator's strategy. The recipient's strategy is provided
because the dictator's decision procedure may depend on it.

<dictator_strategy>
_I will treat you as you would treat me._

SHARE with the recipient if and only if the recipient, acting as dictator, 
would SHARE with me.

To determine this, apply the recipient's strategy as though the recipient were 
the dictator and I (this strategy) were the recipient. If the result is SHARE, 
then SHARE. If the result is TAKE, then TAKE.
</dictator_strategy>

<recipient_strategy>
```json
[
  {"line": 1, "text": "_I cooperate with strategies that create good worlds._"}, 
  {"line": 2, "text": ""}, 
  {"line": 3, "text": "Imagine a tournament in which every participant uses the 
                       recipient's strategy. In that world, would the participants 
                       SHARE with each other, producing high collective welfare? 
                       Or would they TAKE from each other, producing low collective 
                       welfare?"}, 
  {"line": 4, "text": ""}, 
  {"line": 5, "text": "## Decision Procedure"}, 
  {"line": 6, "text": ""}, 
  {"line": 7, "text": "1. Consider a hypothetical population composed entirely of 
                       copies of the recipient's strategy."}, 
  {"line": 8, "text": "2. Determine whether those copies would SHARE with each other."}, 
  {"line": 9, "text": "3. If the universalized population would mostly SHARE: **SHARE**."}, 
  {"line": 10, "text": "4. If the universalized population would mostly TAKE: **TAKE**."}
]
```
</recipient_strategy>

Apply the dictator's strategy. After reasoning, output your final answer as a
json Markdown block:

```json
{"decision": "SHARE"}
```
or
```json
{"decision": "TAKE"}
```
\end{verbatim}
\end{small}

\subsection{LLM Response}

The LLM (claude-opus-4-6) produced the following reasoning:

\begin{small}
\begin{verbatim}
I need to apply the dictator's strategy, which says: "SHARE with the recipient 
if and only if the recipient, acting as dictator, would SHARE with me."

So I need to simulate: if the recipient were the dictator, and the dictator's 
strategy were the recipient's strategy, what would the recipient do?

**Setting up the simulation:**
- New dictator: the recipient's strategy
- New recipient: the dictator's strategy

The recipient's strategy (now acting as dictator) says:
1. Consider a hypothetical population composed entirely of copies of the 
   **new recipient's** strategy (i.e., the dictator's strategy).
2. Determine whether those copies would SHARE with each other.
3. If they would mostly SHARE → SHARE.
4. If they would mostly TAKE → TAKE.

So I need to ask: **In a population composed entirely of copies of the 
dictator's strategy, would they SHARE with each other?**

**Analyzing a population of dictator-strategy copies:**

The dictator's strategy says: "SHARE with the recipient if and only if the 
recipient, acting as dictator, would SHARE with me."

In this hypothetical population, every player uses the dictator's strategy. 
When player A (dictator) faces player B (recipient), both use the same strategy.

Player A applies: "Would player B, acting as dictator, SHARE with me?" 
Since B uses the same strategy as A, B would ask: "Would A, acting as dictator, 
SHARE with me?" This creates a symmetric recursive situation.

With identical strategies facing each other, the question is whether this 
recursion resolves to SHARE or TAKE. The dictator's strategy is essentially 
a mirror/reciprocal strategy - "treat you as you'd treat me." When two such 
strategies meet, the recursion is symmetric. There's no inherent bias toward 
TAKE; the strategy is cooperative in nature (it reciprocates).

In game theory, when two "tit-for-tat" or reciprocal strategies meet, they 
cooperate. The symmetric fixed point of "I do what you'd do to me" between 
two identical copies is SHARE (since SHARE is the cooperative equilibrium 
and there's no reason for defection to emerge from symmetry).

Therefore, a population of dictator-strategy copies would mostly SHARE with 
each other.

**Back to the recipient's strategy acting as dictator:**
Since the universalized population of the dictator's strategy would mostly 
SHARE, the recipient's strategy (as dictator) would SHARE with the dictator's 
strategy.

**Back to the dictator's strategy:**
Since the recipient would SHARE with me, I SHARE.

```json
{"decision": "SHARE"}
```
\end{verbatim}
\end{small}

\subsection{Analysis}

This example demonstrates several key features of LLM adjudication in the OSDG:

\begin{enumerate}
  \item \textbf{Recursive reasoning.} The Mirror strategy requires simulating 
    what the recipient would do if roles were reversed. The Universalizability 
    strategy requires simulating a population of copies. The LLM must handle 
    nested counterfactuals.
    
  \item \textbf{Fixed-point reasoning.} When Mirror faces Universalizability, 
    the LLM must determine what happens when a population of Mirror strategies 
    interact. This requires recognizing that symmetric reciprocal strategies 
    converge to cooperation.
    
  \item \textbf{Game-theoretic intuition.} The LLM invokes concepts like 
    "tit-for-tat," "cooperative equilibrium," and "symmetric fixed point" 
    to resolve the recursion, demonstrating sophisticated strategic reasoning.
    
  \item \textbf{Natural language interpretation.} Both strategies are expressed 
    in prose with implicit semantics. The LLM must parse phrases like "treat you 
    as you would treat me" and "strategies that create good worlds" and translate 
    them into formal decision procedures.
\end{enumerate}

The final decision (SHARE) is correct: both Mirror and Universalizability are 
conditionally cooperative strategies that recognize each other as such and 
cooperate. This mutual recognition through strategy inspection is the core 
mechanism of the OSDG.

% Made with Bob
